%%%%%%%%%%%%%%%%%%%%%%%%%%%%%%%%%%%%%%%%%%%%%%%%%%%%%%%%%%%%%%%%%%%%%%%%%%%%%%%%
%%  Štruktúra práce (kostra). Každý \section je kapitola (nová strana).
%%  Riadky začínajúce % sú poznámky ČO napísať a z čoho - text píšeš sám.
%%  Pozn.: opatrenie STU 1/2024 a poznámka v šablóne - odseky vygenerované AI
%%  do práce nepatria; použitie AI (napr. pri nástrojoch/skriptoch) uveď v
%%  includes/ai_declaration.tex.
%%%%%%%%%%%%%%%%%%%%%%%%%%%%%%%%%%%%%%%%%%%%%%%%%%%%%%%%%%%%%%%%%%%%%%%%%%%%%%%%

%===============================================================================
\section{Analýza problematiky}\label{sec:analyza}
%===============================================================================
% Cieľ kapitoly: teória + súčasný stav, aby čitateľ rozumel kapitolám 2-5.
% Rozsah orientačne 15-20 strán. Každé tvrdenie s citáciou.

\subsection{Legacy systémy a ich testovanie}\label{sec:legacy}
% - definícia legacy kódu (Feathers: "kód bez testov"), prečo je zmena riziková
% - typické vlastnosti: chýbajúca dokumentácia, pevné závislosti, statický stav,
%   dlhé metódy, priamy prístup k DB/súborom/času
% - "seams" (švy) a prečo bez nich nejde jednotkový test
% - charakterizačné testy / golden master: zachytenie súčasného správania
%   pred refaktoringom (aktuálne správanie = očakávané správanie)

\subsection{Jednotkové testovanie v jazyku Java}\label{sec:unittest}
% - čo je jednotkový test vs. integračný, AAA vzor
% - JUnit 5 (Jupiter, parametrizované testy), JUnit 4 (Randoop/EvoSuite)
% - test doubles: stub, mock, fake; Mockito vrátane mockStatic/mockConstruction
% - testovanie JDBC kódu: in-memory databáza H2

\subsection{Hodnotenie kvality testov}\label{sec:metriky-testov}
% - pokrytie riadkov a vetiev (JaCoCo), prečo pokrytie nestačí
% - mutačné testovanie: mutant, zabitý/prežitý mutant, mutačné skóre (PIT),
%   vzorec skóre = detegované / (všetky - nevalidné)
% - problém orákula (test oracle problem): kto rozhodne, čo je správny výsledok
% - spoľahlivosť testov: flaky testy, závislosť od času/náhody/poradia

\subsection{Testovateľnosť a štrukturálne vlastnosti kódu}\label{sec:testovatelnost}
% - definícia testovateľnosti
% - cyklomatická komplexita (McCabe), kognitívna komplexita (SonarSource)
% - väzby a závislosti (coupling, fan-out), hard-wired závislosti, statický stav
% - väzba na zadanie: "vplyv štrukturálnych vlastností na testovateľnosť"

\subsection{Klasické nástroje na automatické generovanie testov}\label{sec:klasicke}
% - náhodné generovanie s regresnými orákulami: Randoop
% - search-based (genetický algoritmus, kritériá pokrytia): EvoSuite
% - ich obmedzenia: nečitateľné testy, orákulum = súčasný výstup,
%   EvoSuite len po bajtkód Java 11 (dôležité pre prípadovú štúdiu)

\subsection{Veľké jazykové modely pre generovanie testov}\label{sec:llm}
% - LLM v skratke: transformer, tokeny, kontextové okno, teplota, kvantizácia
% - modely pre kód; lokálny beh (Ollama) vs. cloudové API (súkromie kódu!)
% - prompt engineering: rola, pravidlá, kontext, few-shot
% - spätnoväzbová slučka (kompilátor + výsledky testov -> oprava)
% - súvisiace práce: TestPilot, ChatUniTest, CodaMosa, CoverUp, Meta TestGen-LLM,
%   prehľadové štúdie o LLM v testovaní - čo merali a čo im chýba
% - medzera, ktorú práca vypĺňa: legacy kód so zlou testovateľnosťou,
%   kontrolovaný benchmark + mutačné skóre + lokálne modely

\subsection{Zhrnutie analýzy a výskumné otázky}\label{sec:rq}
% Sformuluj výskumné otázky (budú sa vracať v kapitole Výsledky):
% RQ1 Úspešnosť: koľko vygenerovaných testov sa skompiluje a prejde?
% RQ2 Kvalita: aké pokrytie a mutačné skóre dosahujú LLM v porovnaní
%     s klasickými nástrojmi a golden masterom?
% RQ3 Spoľahlivosť: aké stabilné sú výsledky medzi iteráciami (rozptyl)?
% RQ4 Testovateľnosť: ako štrukturálne vlastnosti triedy ovplyvňujú výsledok?
% (RQ5 Refaktoring: dokáže AI upraviť kód, aby bol lepšie testovateľný?
%  - je v abstrakte, experiment ešte NIE JE spravený -> spraviť alebo vyhodiť z abstraktu)

%===============================================================================
\section{Návrh metodiky}\label{sec:metodika}
%===============================================================================

\subsection{Benchmark legacy tried}\label{sec:benchmark}
% - prečo vlastný kontrolovaný benchmark (každá trieda = jedna vlastnosť)
% - tabuľka B01-B15: trieda | skúmaná vlastnosť | riadky | cyklomatická
%   komplexita | závislosti   (údaje: zdrojáky v benchmark-core/src/main/java)
% - core modul kompilovaný pre Java 8 (rovnaký bajtkód pre všetky nástroje),
%   B09 moderná Java v samostatnom module

\subsection{Porovnávané prístupy}\label{sec:pristupy}
% B0 bez testov | B1 golden master | B2 Randoop, EvoSuite | B3 naivný prompt |
% M navrhnutá metodika | ručne písané referenčné testy (horná hranica)
% Pri každom: čo to je a prečo je v porovnaní.

\subsection{Navrhnutá metodika generovania}\label{sec:metodika-m}
% - naivný prompt vs. plný prompt (rola, pravidlá, kontext použitých tried)
% - odstránenie komentárov zo zdrojového kódu (férovosť voči Randoop/EvoSuite)
% - opravná slučka: max. 3 kolá, ponechá sa najlepšia verzia
% - obrázok: diagram toku generovanie -> kontrola -> oprava

\subsection{Metriky}\label{sec:metriky}
% podiel skompilovaných súborov, pass rate, pokrytie riadkov/vetiev,
% mutačné skóre, čas generovania, počet tokenov; definuj každú presne

\subsection{Návrh experimentu}\label{sec:experiment}
% 5 iterácií, seed = 41 + iterácia, teplota 0.2, rozpočet 60 s/trieda,
% Randoop limit 500 testov, kontext 16384 tokenov
% modely: Qwen3.6-35B-A3B, Qwen3-Coder-30B-A3B, phi4-mini (parametre, kvantizácia)
% HW a SW: MacBook (CPU, RAM), JDK 25, JaCoCo 0.8.14, PIT 1.25.8 ...
% (všetko je v results/runs.csv - stĺpce os, cpu_count, ram_gb, java_version)

%===============================================================================
\section{Implementácia meracieho prostredia}\label{sec:implementacia}
%===============================================================================

\subsection{Architektúra}\label{sec:architektura}
% obrázok pipeline: generovanie -> kompilácia -> testy + JaCoCo -> PIT -> CSV -> report

\subsection{Generovanie testov}\label{sec:impl-generovanie}
% skripty generate_randoop.sh, generate_evosuite.sh, generate_llm.py,
% approve_golden_master.sh; Ollama, streamovanie, metadáta generovania

\subsection{Meranie}\label{sec:impl-meranie}
% run_iteration.py: odstraňovanie neskompilovateľných súborov, izolácia
% tried, ktoré zhodia JVM, PIT s časovým limitom, výstupné CSV, HTML report

\subsection{Reprodukovateľnosť}\label{sec:reprodukovatelnost}
% generation_meta.txt, verzie nástrojov, seedy, upozornenie na zmiešané nastavenia

\subsection{Prekonané problémy}\label{sec:problemy}
% EvoSuite na JDK 25 -> JDK 8; Randoop 28 tisíc testov -> limit;
% nekonečný výstup malého modelu -> limit tokenov; mockStatic(Math) zhodil JVM;
% opravná slučka zhoršovala výsledok -> ponechanie najlepšej verzie

%===============================================================================
\section{Výsledky}\label{sec:vysledky}
%===============================================================================
% Údaje: legacy-test-benchmark/results/runs.csv, classes.csv, report.html
% Pri každom čísle priemer ± smerodajná odchýlka z 5 iterácií.

\subsection{Prehľad výsledkov}\label{sec:prehlad}
% súhrnná tabuľka: prístup | pass rate | pokrytie riadkov | vetiev | mutačné skóre | čas

\subsection{Úspešnosť generovania (RQ1)}\label{sec:rq1}
\subsection{Pokrytie a mutačné skóre (RQ2)}\label{sec:rq2}
\subsection{Stabilita medzi iteráciami (RQ3)}\label{sec:rq3}
\subsection{Vplyv štrukturálnych vlastností (RQ4)}\label{sec:rq4}
% graf: mutačné skóre po triedach; príklady: InvoiceProcessor (God metóda),
% CustomerRepository (JDBC) - koľko opravných kôl, prečo zlyhali testy
\subsection{Prípadová štúdia na reálnom projekte}\label{sec:pripadova}
% LLM vs. Randoop (EvoSuite nepodporuje novšiu Javu), 3-5 -> 10-15 tried,
% ručná kontrola vzorky padajúcich testov: chybné očakávanie vs. nájdená chyba

%===============================================================================
\section{Diskusia}\label{sec:diskusia}
%===============================================================================

\subsection{Interpretácia výsledkov}\label{sec:interpretacia}
% kedy sa LLM oplatí, kedy golden master, kedy klasický nástroj

\subsection{Odporúčania pre prax}\label{sec:odporucania}
% lokálne modely kvôli súkromiu kódu, opravná slučka je nutná, kontrola človekom

\subsection{Hrozby validity}\label{sec:validita}
% vnútorná (seedy, nedeterminizmus LLM, HW), vonkajšia (malý benchmark,
% jeden reálny projekt), konštrukčná (mutačné skóre ako proxy kvality),
% únik dát (modely mohli vidieť podobný kód pri trénovaní)
